\documentclass[a4paper, oneside]{article}
\usepackage{graphicx}
\usepackage{amsthm}
\usepackage{amsmath}
\usepackage{amssymb}
\usepackage[a4paper,
            bindingoffset=0.2in,
            left=2cm,
            right=2cm,
            top=2cm,
            bottom=2cm,
            footskip=.25in]{geometry}
\usepackage[italian]{babel}
\usepackage{pgfplots}
\usepackage{tabularx}
\usepackage{tikz}
\usepackage{wrapfig}
\usepackage{color}
\usepackage[d]{esvect}
\usepackage{circuitikz}
\definecolor{page}{rgb}{0.129,0.157,0.212}
\pagecolor{page}
\color{white}
\graphicspath{ {./images/} }
\usetikzlibrary{shapes.geometric}
\usetikzlibrary{datavisualization}
\usetikzlibrary{datavisualization.formats.functions}
\usetikzlibrary{patterns}
\pgfplotsset{width=10cm,compat=1.18}

\title{Appunti di Lab II}
\author{Tommaso Miliani}
\date{09-03-26}

\begin{document}
\newtheoremstyle{theoremEnv}
                {}          % Space above
                {}          % Space below
                {\slshape}  % Body font
                {}          % Indent amount
                {\bfseries} % Head font
                {.}         % Punctuation after head
                {\newline}  % Space after theorem head
                {}          % Theorem head spec
\theoremstyle{theoremEnv}

\newtheorem{definition}{Definizione}[section]
\newtheorem{theorem}{Teorema}[section]
\newtheorem{lemma}{Proposizione}[section]
\newtheorem{observation}{Osservazione}[section]
\newtheorem{corollary}{Corollario}[theorem]
\newtheorem{example}{Esempio}[section]
\newtheorem{remark}{Enunciato}[section]

\maketitle

\section{Esercizio}
\begin{wrapfigure}{r}{0.4\textwidth}
    \centering
    \caption{}
    \begin{circuitikz}
        \draw
        (0,0) to [short, *-] (6,0)
        to [V, l_=$\mathrm{j}{\omega}_m \underline{\psi}^s_R$] (6,2) 
        to [R, l_=$R_R$] (6,4) 
        to [short, i_=$\underline{i}^s_R$] (5,4) 
        (0,0) to [open, v^>=$\underline{u}^s_s$] (0,4) 
        to [short, *- ,i=$\underline{i}^s_s$] (1,4) 
        to [R, l=$R_s$] (3,4)
        to [L, l=$L_{\sigma}$] (5,4) 
        to [short, i_=$\underline{i}^s_M$] (5,3) 
        to [L, l_=$L_M$] (5,0); 
    \end{circuitikz} 
\end{wrapfigure}
Il metodo delle correnti di maglia permettono di risolvere completamente il circuitop
con un set minimo di equazioni per ciascuna maglia. Per ogni maglia si definisce una corrente 
di maglia con un certo senso in modo da trovare una matrice $n \times n$ come sistema lineare.
\begin{gather*}
    \left\{\begin{array}{l}
    \epsilon_1 + \epsilon_2 = I_1 R_1 + (I_1 + I_2) R_2  \\
    -\epsilon_2 = I_2 R_3 + I_2 R_2 - I_1 R_2 
    \end{array}\right. \qquad I_2 = i_3
\end{gather*}
Si può anche riscrivere, in forma matriciale:
\begin{gather*}
    \left\{\begin{array}{l}
        \epsilon_1 + \epsilon_2 = (R_1 + R_2) I_1 + (-R_2) I_2 \\
        -\epsilon_2 = (-R_2)I_1 + (R_2 + R_3) I_2
    \end{array}\right. \ \Longrightarrow \ \\
    \begin{pmatrix}
        R_1 + R_2 & -R_2 \\
        -R_2 & R_2 + R_3
    \end{pmatrix} \begin{pmatrix}
        I_1 \\
        I_2
    \end{pmatrix} = \begin{pmatrix}
        \epsilon_1 + \epsilon_2  \\
        - \epsilon_2 
    \end{pmatrix} \\
    \ \Longrightarrow \ \epsilon_{i, 1} = R_{i, j} I_{j, 1}
\end{gather*}
L'ultimo termine è il prodotto tra matrici semplificato. Questo metodo è generale e 
la soluzione utilizza il metodo di Cramer:
\begin{gather*}
    I_2 = i_3 = 
\end{gather*}

PEzzo mancante

\section{Teorema di Thevenin}
Il \textbf{teorema di Thevenin} mi dice che ogni qualsiavoglia 
rete è equivalente ad un genratore di $fem$ equivalente $\epsilon_{TH}$ e 
una resistenza equivalente $R_{TH}$: dunque qualsiasi rete è schematizzabile 
come se fosse solamente un generatore reale con resistenza interna. 
Si introduce un Esempio della trasformazione triangolo stella. 
Si dimostra ora il teorema: \\
La resistenza totale è quella
equivalente se tutti i generatoria sono cortocircuitati. $\epsilon_{TH}$ + la differenza di potenziale 
tra $A$ e $B$ se $i = 0$. Se si cortocircuita tutto il circuito si ottiene una corrente
$I_{cc}$ di cortocircuito. Allora 
\begin{align}
    R_{TH} = \frac{\epsilon_{TH}}{I_{cc}}
\end{align}
Si costruisce una rete aggiuntiva (ossia a cui si aggiunge un generatore reale), 
la soluzione per $I_n$ è il rapporto tra i determinanti, dunque si deve sostituire,
all'ultima colonna, la colonna dei termini noti:
ali


\end{document}